\chapter{Grafikus felület specifikációja}

\section{Fordítási és futtatási útmutató}
\comment{A feltöltött program fordításával és futtatásával kapcsolatos útmutatás. Ennek tartalmaznia kell leltárszerűen az egyes fájlok pontos nevét, méretét byte-ban, keletkezési idejét, valamint azt, hogy a fájlban mi került megvalósításra.}

\subsection{Fájllista}

\begin{fajllista}
    \fajl
    {Demo.java} % Kezdet
    {353 byte} % Méret
    {2020.03.26~21:05~} % Időpont
    {Demo programosztály} % Leírás
\end{fajllista}

\subsection{Fordítás}
\begin{verbatim}
A projekt fő mappájában kell ezt a parancsot egyben futtatni(NAGYON FONTOS: a parancsot egyben kell futtatni, különben nem fog működni! Másold ki txt-be és töröld ki a sortöréseket!):
javac -d bin prototipusVegleges/src/jarmuvek/*java 
prototipusVegleges/src/jatekosok/*.java 
prototipusVegleges/src/motor/*.java 
prototipusVegleges/src/takaritofejek/*.java 
prototipusVegleges/src/terkep/*.java 
prototipusVegleges/src/gui/*.java 
prototipusVegleges/src/gui/controller/*.java 
prototipusVegleges/src/gui/view/*.java
\end{verbatim}

\subsection{Futtatás}
\begin{verbatim}
cd bin
java gui.Main
\end{verbatim}

\clearpage
\section{Értékelés}
\begin{ertekeles}
    \ertekelestag{Görgényi András Levente}{ GSQEM0 }{20\%}
    \ertekelestag{Dulcz Bence}{ Y9LPT9 }{20\%}
    \ertekelestag{Varga Gergely}{ IBXDS0 }{20\%}
    \ertekelestag{Nagy Patrik}{ C3IJYQ }{20\%}
    \ertekelestag{Tóth Dávid Pál}{ H9JYWA }{20\%}
\end{ertekeles}

